\documentclass[12pt,a4paper]{article}

\usepackage{fontspec}
\setmainfont{Liberation Serif}
\setsansfont{Liberation Sans}
\setmonofont{Liberation Mono}
\linespread{1.0}
\fontsize{14pt}{17pt}\selectfont

\usepackage{geometry}
\geometry{left=3cm,right=1.5cm,top=2cm,bottom=2cm}

\usepackage{indentfirst}
\setlength{\parindent}{1.25cm}

\usepackage{array}
\usepackage{tabularx}

\usepackage{fancyhdr}
\pagestyle{fancy}
\fancyhf{}
\fancyfoot[C]{\thepage}
\renewcommand{\headrulewidth}{0pt}

\begin{document}

% Титульная страница
\thispagestyle{empty}
\begin{center}
ГОСУДАРСТВЕННОЕ УЧРЕЖДЕНИЕ ОБРАЗОВАНИЯ\\
<<МИНСКИЙ ГОСУДАРСТВЕННЫЙ КОЛЛЕДЖ\\
ЭЛЕКТРОНИКИ>>
\end{center}

\vspace{3cm}

\begin{center}
{\bfseries ОТЧЁТ\\
ПО АВТОМАТИЗИРОВАННЫМ ТЕСТАМ\\
ПРОЕКТА <<GoldPC>>}
\end{center}

\vspace{2cm}

\begin{center}
2026
\end{center}

\vfill

\newpage

\section*{СОДЕРЖАНИЕ}

1.~Назначение автоматизированных тестов\\
2.~Средства автоматизации\\
3.~Структура тестовой инфраструктуры\\
4.~Результаты автоматизированных тестов\\
5.~Выводы

\newpage

\section{1. Назначение автоматизированных тестов}

Автоматизированные тесты в проекте <<GoldPC>> обеспечивают корректность бизнес-логики микросервисов, предотвращают регрессии при изменении кода и валидируют архитектурные ограничения.

\section{2. Средства автоматизации}

Для написания и запуска тестов используются: xUnit (фреймворк тестирования), FluentAssertions (утверждения), Moq (моки), AutoFixture (тестовые данные), EF Core InMemory (in-memory БД), WebApplicationFactory (HTTP-тестирование), PactNet (контрактные тесты), NetArchTest.Rules (архитектурные тесты), Bogus (Fakers).

\section{3. Структура тестовой инфраструктуры}

Тестовая инфраструктура организована в виде отдельных проектов: AuthService.Tests, GoldPC.UnitTests, OrdersService.Tests, ServicesService.Tests, WarrantyService.Tests, PCBuilderService.Tests, ArchitectureTests, GoldPC.IntegrationTests, GoldPC.ContractTests.

\section{4. Результаты автоматизированных тестов}

В таблице представлены 7 основных модульных тестов.

\begin{center}
\begin{tabularx}{\textwidth}{|c|X|c|}
\hline
№ & Тест & Результат \\
\hline
1 & Register\_ValidRequest\_ShouldCreateUser & Пройден \\
\hline
2 & GetProductById\_ReturnsProductDetail & Пройден \\
\hline
3 & CreateOrder\_ValidRequest\_ShouldCreateOrder & Пройден \\
\hline
4 & CreateServiceRequest\_ValidRequest & Пройден \\
\hline
5 & CreateWarrantyClaim\_ValidRequest & Пройден \\
\hline
6 & CheckCompatibility\_IncompatibleSocket & Пройден \\
\hline
7 & ServicesShouldHaveInterfaces & Пройден \\
\hline
\end{tabularx}
\end{center}

Всего тестов: 7. Пройдено: 7. Не пройдено: 0.

\section{5. Выводы}

Автоматизированная тестовая инфраструктура проекта <<GoldPC>> обеспечивает высокий уровень доверия к корректности бэкенд-системы. Все 7 ключевых модульных тестов проходят успешно, что подтверждает надёжность реализации основной бизнес-логики микросервисов.

\end{document}
